\subsection{Item Mode Panel}
\label{sec:geo_item_mode}

The Item Mode panel sits in the bottom half of the 'Layers' tab splitter and acts as a per-target-report item browser. It groups the loaded target reports by a chosen attribute and lets the user pick individual items, jump to them on the map and run item-specific actions. \\\\

\begin{figure}[H]
  \center
    \fbox{\parbox{0.8\textwidth}{\vspace{1.5cm}\centering \textcolor{red}{\textbf{[FIGURE TODO: Layers tab splitter with the Layers tree on top and the Item Mode panel below]}}\vspace{1.5cm}}}
  \caption{Geographic View: Item Mode panel}
\end{figure}

At the top of the panel an 'Item Mode' combo box selects how items in the list are grouped. The following groupings are available:

\begin{itemize}
\item None
\item Aircraft Address
\item Aircraft Identification
\item Track Number
\item Mode 3/A Code
\item UTN
\end{itemize}
\ \\

\begin{figure}[H]
  \center
    \fbox{\parbox{0.8\textwidth}{\vspace{1.5cm}\centering \textcolor{red}{\textbf{[FIGURE TODO: 'Item Mode' combo expanded showing all six grouping options]}}\vspace{1.5cm}}}
  \caption{Geographic View: Item Mode combo}
\end{figure}

Below the combo, a tree view lists the items resulting from the selected grouping. Interaction with an item is:

\begin{itemize}
\item Double-click on an item: zooms the map to that item
\item Right-click on an item: opens a context menu with item-specific actions (zoom, hide / show and other actions provided by the underlying geometry item provider)
\end{itemize}
\ \\

\begin{figure}[H]
  \center
    \fbox{\parbox{0.8\textwidth}{\vspace{1.5cm}\centering \textcolor{red}{\textbf{[FIGURE TODO: right-click context menu on an item in the Item Mode panel showing the available actions]}}\vspace{1.5cm}}}
  \caption{Geographic View: Item Mode context menu}
\end{figure}

Typical use cases include finding a specific target by Aircraft Identification or Aircraft Address, jumping to a UTN's path, or focusing on a single Mode 3/A code. \\

Please \textbf{note} that the Item Mode panel is a per-target-report item browser; it is not an additional Color Mode level (see \nameref{sec:color_mode}). \\

Please also \textbf{note} that for very large numbers of items the list can become slow; no hard limit is enforced. If interactive responsiveness is more important than browsing every item, a coarser grouping (e.g. 'None' or 'UTN') can be selected.
